\documentclass[10pt,a4paper]{article}
\usepackage{amsmath,amssymb,amsthm,graphicx}

\usepackage{amsmath,amssymb,amsfonts,amsthm,mathtools,stmaryrd,xspace,enumitem}
\newenvironment{authorclaim}[1]{\par\medskip\noindent\textbf{#1.}\ \ignorespaces}{\par\medskip}
\newcommand{\oneset}[1]{\left\{\mathinner{#1}\right\}}
\newcommand{\ccinterval}[2]{[{#1}..{#2}]}
\newcommand{\abs}[1]{\left|\mathinner{#1}\right|}
\newcommand{\N}{\ensuremath{\mathbb N}}
\newcommand{\eps}{\varepsilon}
\newcommand{\Oh}{\mathcal O}
\newcommand{\Nc}[1]{\ensuremath{\mathsf{NC}^{#1}}\xspace}
\newcommand{\Ac}[1]{\ensuremath{\mathsf{AC}^{#1}}\xspace}
\newcommand{\DLINTIME}{\ensuremath{\mathsf{DLINTIME}}\xspace}
\newcommand{\WP}{\textsc{WP}\xspace}
\newcommand{\DTIME}{\ensuremath{\mathsf{DTIME}}\xspace}
\newcommand{\nand}{\ensuremath{\mathsf{nand}}}
\newbox\pairbox
\def\pair<#1>{{\mathsurround=0pt\setbox\pairbox\hbox{$\left\langle#1\right\rangle$}\left\langle\kern-0.35\ht\pairbox\copy\pairbox\kern-0.35\ht\pairbox\right\rangle}}
\AtBeginDocument{\author{\parbox{.95\textwidth}{\centering Laurent Bartholdi, Michael Figelius\\Markus Lohrey, Armin Wei\ss}}}
\allowdisplaybreaks
\newtheorem{definition}{Definition}
\title{Nonuniform NC1 projection hardness of SENS group word problems via nested-commutator programs}
\author{Laurent Bartholdi, Michael Figelius, Markus Lohrey, Armin Weiß}
\date{doi:10.4230/LIPIcs.CCC.2020.29; CC BY 3.0 Unported}
\begin{document}
\maketitle
\noindent\textit{Editorial scope.} One independent nonuniform NC1 projection mechanism with full SENS witnesses grouped. Does not claim complete uniform ALOGTIME theorem or PSPACE corollaries.\par\medskip
\noindent\textit{Allowed standard prerequisites (editorial).}\begin{itemize}
\item Group operations, balanced Boolean NAND formulas and standard machine/circuit normal form. Formal standard finite symmetric generating set includes identity for padding. Original references[3,12,17,47] retained as standard prerequisites; full new nested-commutator program and both explicit group constructions retained.
\end{itemize}\medskip
\section*{Formal non-uniform conclusion}
Since the proof for the uniform case is not difficult but tedious, in this sketch we primarily focus on the non-uniform version, which only gives us hardness via (non-uniform) $\mathrm{AC}^0$-reductions instead of DLOGTIME-reductions.
This gives us a non-uniform $\mathrm{AC}^0$-reduction (more precisely, a projection reduction) of $L(M)$ to $\mathrm{WP}(G)$.

\section*{Standard generating sets and word problems}
Recall that a group $G$ with identity element $1$ is finitely generated (f.g.\ for short) if there is a finite set $\Sigma\subseteq G$ such that every element of $G$ can be written as a product of elements of $\Sigma$; this product can be formally written as a word from $\Sigma^*$. For technical reasons we assume that $1\in\Sigma$ (which is needed for padding reasons) and that for every $a\in\Sigma$ also the inverse $a^{-1}$ belongs to $\Sigma$; such a generating set $\Sigma$ is called standard. We have a natural involution on $\Sigma^*$ defined by $(a_1\cdots a_n)^{-1}=a_n^{-1}\cdots a_1^{-1}$ for $a_i\in\Sigma$ (which is the same as forming inverses in the group). For words $u,v\in\Sigma^*$ we write $u=_Gv$ if $u$ and $v$ are equal in $G$; sometimes we just say $u=v$ in $G$. The word problem for $G$, $\mathrm{WP}(G)$ for short, is the question whether $u=_G1$ holds for a given word $u\in\Sigma^*$. In this formulation the word problem depends on the generating set $\Sigma$, but it is well-known that the complexity/decidability status of the word problem does not depend on $\Sigma$.
\section*{SENS conventions (Definition 1 and Section 4)}
For group elements $g,h\in G$ or words $g,h\in\Sigma^*$ we write $g^h$ for the conjugate $h^{-1}gh$ and $[h,g]$ for the commutator $h^{-1}g^{-1}hg$. We call $g$ a $d$-fold nested commutator, if $d=0$ or $g=[h_1,h_2]$ for $(d-1)$-fold nested commutators $h_1,h_2$. A subquotient of a group $G$ is a quotient of a subgroup of $G$.
\begin{definition}\label{def:SENS}		
  We call a group $G$ with the finite standard generating set $\Sigma$
  \emph{strongly efficiently non-solvable} \emph{(SENS)} if for every
  $d \in \N$ there is a collection of $2^{d+1}-1$ elements
  $g_{d,v} \in \Sigma^*$ for $v \in \oneset{0,1}^{\leq d}$ such that
  \begin{enumerate}[label=(\alph*)]
  \item\label{SENSa} there is some constant $\mu \in \N$ with $\abs{g_{d,v}} = 2^{\mu d}$ for all $v \in \oneset{0,1}^{d}$,
  \item\label{SENSb} $g_{d,v} =	\bigl[ g_{d,v0},\,  g_{d,v1} \bigr]$ for all $v \in \oneset{0,1}^{< d}$ (here we take the commutator of words),
  \item\label{SENSc} $g_{d,\eps} \neq 1$ in $G$.
  \end{enumerate}
   The group $G$ is called \emph{uniformly strongly efficiently
    non-solvable} if, moreover,
  \begin{enumerate}[label=(\alph*)]
    \setcounter{enumi}{3}
  \item\label{SENSu} given $v\in \oneset{0,1}^d$, a number $i$ encoded in binary with $\mu d$ bits, and $a \in \Sigma$ one can decide in $\DLINTIME$ whether the
  	$i$-th letter of $g_{d,v}$ is $a$.
  \end{enumerate} 
\end{definition}
Here are some simple observations:
\begin{itemize}
\item A strongly efficiently non-solvable group clearly cannot
be solvable, so the above terminology makes sense.
\item If one can find suitable $g_{d,v}$ of length at most $2^{\mu d}$,  then these words can always be padded to length $2^{\mu d}$ thanks to the padding letter $1$.
\item  It suffices to specify $g_{d,v}$ for $v \in \oneset{0,1}^{d}$;
  the other $g_{d,v}$ are then defined by Condition~(b).
 \item We have $\abs{g_{d,v}} = 2^{\mu d+ 2(d-\abs{v})}$ for all $v \in \oneset{0,1}^{\leq d}$. Thus, all $g_{d,v}$ have length $2^{\Oh(d)}$.
 \item Equivalently to Condition~(d), we can require that given $v\in \oneset{0,1}^d$ and a binary encoded number $i$ with $\mu d$ bits, one can compute the $i$-th letter of $g_{d,v}$ in $\DLINTIME$.
\end{itemize}
Henceforth, whenever $d$ is clear, we simply write $g_v$ instead of $g_{d,v}$.


\section*{Necessary group context (Section 2)}
Richard Thompson introduced three finitely presented groups $F<T<V$ acting on the unit-interval, the unit-circle and the Cantor set, respectively. A standard reference for Thompson's groups is [12]. The group $F$ consists of all homeomorphisms of the unit interval that are piecewise affine, with slopes a power of $2$ and dyadic breakpoints. Famously, $F$ has a finite presentation with two generators:
\[
F=\langle x_0,x_1\mid[x_0x_1^{-1},x_0^{-1}x_1x_0],[x_0x_1^{-1},x_0^{-2}x_1x_0^2]\rangle .
\]
Very convenient is also the following infinite presentation:
\[
F=\langle x_0,x_1,x_2,\ldots\mid x_k^{x_i}=x_{k+1}\ (i<k)\rangle .
\]
The group $F$ is orderable (so in particular torsion-free), its derived subgroup $[F,F]$ is simple and the center of $F$ is trivial.
Let $X$ be a finite set. The free monoid $X^*$ serves as the vertex set of a regular rooted tree with an edge between $v$ and $vx$ for all $v\in X^*$ and all $x\in X$. The group $W$ of automorphisms of this tree naturally acts on the set $X$ of level-1 vertices, and permutes the subtrees hanging from them. Exploiting the bijection $X^+=X^*\times X$, we thus have an isomorphism
\[
\varphi:W\to W\wr\mathrm{Sym}(X)=W^X\rtimes\mathrm{Sym}(X),
\]
mapping $g\in W$ to elements $f\in W^X$ and $\pi\in\mathrm{Sym}(X)$ as follows: $\pi$ is the restriction of $g$ to $X\subseteq X^*$, and $f$ is uniquely defined by $(xv)^g=x^\pi v^{f(x)}$. We always write $g@x$ for $f(x)$ and call it the state (or coordinate) of $g$ at $x$. If $X=\{0,\ldots,k\}$ we write $g=\langle\!\langle g@0,\ldots,g@k\rangle\!\rangle\pi$.
There exist important examples of f.g.\ self-similar weakly branched groups, notably the Grigorchuk group $G$, see [17]. It may be described as a self-similar group in the following manner: it is a group generated by $\{a,b,c,d\}$, and acts on the rooted tree $X^*$ for $X=\{0,1\}$. The action, and therefore the whole group, are defined by the restriction of $\varphi$ to $G$'s generators: $\varphi(a)=(0,1)$, $\varphi(b)=\langle\!\langle a,c\rangle\!\rangle$, $\varphi(c)=\langle\!\langle a,d\rangle\!\rangle$, and $\varphi(d)=\langle\!\langle1,b\rangle\!\rangle$, where we use the notation $(0,1)$ for the non-trivial element of $\mathrm{Sym}(X)$ (that permutes $0$ and $1$) and $\langle\!\langle w_0,w_1\rangle\!\rangle$ for a tuple in $G^{\{0,1\}}\cong G\times G$.

\section*{The complete formal non-uniform construction (Section 4)}
The proof of Theorem 2 essentially follows Barrington's proof that the word problem of finite non-solvable groups is ALOGTIME-hard [3]. The entire proof can be found in the full version [5], where we also state a non-uniform version of Theorem 2. Since the proof for the uniform case is not difficult but tedious, in this sketch we primarily focus on the non-uniform version, which only gives us hardness via (non-uniform) $\mathrm{AC}^0$-reductions instead of DLOGTIME-reductions.
Like in Barrington's proof, we start with an ALOGTIME-machine (resp.\ $\mathrm{NC}^1$-circuit) and construct a family of so-called $G$-programs. Since we are dealing with finitely generated, but infinite groups, we have to adapt the definition of $G$-programs slightly.
Fix a finite standard generating set $\Sigma$ of $G$. A $G$-program $P$ of length $m$ and input length $n$ is a sequence of instructions $\langle i_j,b_j,c_j\rangle$ for $0\le j\le m-1$ where $i_j\in[1..n]$ and $b_j,c_j\in\Sigma$. On input of a word $x=x_1\cdots x_n\in\{0,1\}^*$, an instruction $\langle i_j,b_j,c_j\rangle$ evaluates to $b_j$ if $x_{i_j}=1$ and to $c_j$ otherwise. The evaluation of a $G$-program is the product (in the specified order) of the evaluations of its instructions, and is denoted with $P[x]\in\Sigma^*$.
Let $M$ be an ALOGTIME-machine in input normal form [47, Lemma 2.41], i.e., every computation path queries at most one input bit and $M$ halts immediately after the query. For every input size $n$, the computation tree of $M$ translates immediately into a Boolean circuit of depth $d\in O(\log n)$. Moreover, $M$ can be normalized such that this circuit is a fully balanced binary tree meaning that the gates of the circuit are indexed by the set $\{0,1\}^{\le d}$, where $\{0,1\}^{<d}$ are the inner gates (where $\varepsilon$ is the output gate, which counts here as an inner gate) and $\{0,1\}^d$ are the leaves (input gates). We can assume that all inner gates are nand-gates (where the Boolean function $\mathsf{nand}:\{0,1\}^2\to\{0,1\}$ is defined by $\mathsf{nand}(0,0)=\mathsf{nand}(0,1)=\mathsf{nand}(1,0)=1$ and $\mathsf{nand}(1,1)=0$) and each leaf is labelled by a possibly negated input variable or constant via an input mapping $q_n:\{0,1\}^d\to[1..n]\times\{0,1\}\times\{0,1\}$. The meaning of this mapping is as follows: if the input to the circuit is the bit string $x_1x_2\cdots x_n$ and $q_n(v)=\langle i,a,b\rangle$, then the input gate $v\in\{0,1\}^d$ evaluates to $a$ (resp., $b$) if $x_i=1$ (resp., $x_i=0$).
The family of circuits obtained this way can be shown to be DLOGTIME-uniform in an even stronger sense than the usual definition (see e.g.\ [47]). For the sake of a simpler description, we fix the input length $n$ and write $C$ for the $n$-input circuit of depth $d\in O(\log n)$. W.l.o.g.\ for every $x\in\{0,1\}^n$, we have $x\in L(M)$ if and only if the output gate of $C$ evaluates to $0$ on input $x$.
\begin{proof}
  For each gate $v \in \oneset{0,1}^{\leq d}$ let
  $g_v = g_{d,v}$ as in Definition 1. We construct two $G$-programs $P_v$
  and $P_v^{-1}$ (both of input length $n$) such that for every input $x \in \{0,1\}^n$
  ($x$ is taken as the input for $C$,  $P_v$, and $P_v^{-1}$) we have
  \begin{align}\label{eq:pvdef}\tag{2}
    P_v[x] &=_G \begin{cases}
      g_v &\text{if } v \text{ evaluates to } 1,\\
      1 &\text{if } v \text{ evaluates to } 0, 		
    \end{cases} 
  \end{align}
  and $P_v^{-1}[x] = P_v[x]^{-1}$ in $G$.  Notice that we have
  $g_v P_v^{-1}[x] = g_v$ if $v$ evaluates to $0$ and $g_v P_v^{-1}[x] = 1$,
  otherwise. Thus, $g_v P_v^{-1}$ is a $G$-program for the ``negation''
  of $P_v$. Moreover, by Equation (2), $P_\eps$ evaluates to $1$ on input $x$ if and only if the output gate evaluates to $0$ which by our assumption was the case if and only if  $x \in L$.
	
  The construction of the $P_v$ and $P_v^{-1}$ is straightforward:
  For an input gate $v \in \oneset{0,1}^{d}$ we simply define $P_v$ to be a $G$-program evaluating to $g_v$ or $1$~-- in which case it evaluates to which element depends on $q_n(v)$. %
  More precisely, write $g_v = a_1\cdots a_m$ with $a_i \in \Sigma$. If $q_n(v) = \langle i, a,b\rangle$ for $i \in \ccinterval{1}{n}$ and $a,b \in \oneset{0,1}$, we set 
  $P_v =  \langle i, a_1^a,a_1^b\rangle \cdots  \langle i, a_m^a,a_m^b\rangle$ and $P_v^{-1} =  \langle i, a^{-a}_m,a^{-b}_m\rangle \cdots  \langle i, a^{-a}_1,a^{-b}_1\rangle$.
  For a {\sf nand}-gate $v$ with
  inputs from $v0$ and $v1$, we define
  \begin{align*}
    P_v &= g_v [P_{v1}, P_{v0}]=g_v  P_{v1}^{-1}  P_{v0}^{-1} P_{v1} P_{v0}, \\
    P_v^{-1} &=  [P_{v0}, P_{v1}] g_v^{-1}= P_{v0}^{-1} P_{v1}^{-1}  P_{v1} P_{v1} g_v^{-1},
  \end{align*}
  where the $g_v$  and $g_v^{-1}$ represent constant
  $G$-programs evaluating to $g_v$ and $g_v^{-1}$, respectively, irrespective of the actual input
  (such constant $G$-programs consist of triples of the form $\langle 1, a,a\rangle$ for $a \in \Sigma$). These
  constant $G$-programs are defined via the commutator identities
  $g_v = \bigl[ g_{v0},\, g_{v1} \bigr]$ for $v \in \oneset{0,1}^{< d}$ in
  Definition 1.

  Clearly, by induction we have $ P_v[x]^{-1} = P_v^{-1}[x]$ in $G$ (for every input $x$). Let us
  show that Equation (2) holds: For an input gate
  $v \in \oneset{0,1}^{d}$, Equation (2) holds by definition. Now, let
  $v \in \oneset{0,1}^{< d}$. Then, by induction, we have the following equalities in $G$:
  \begin{align*}
    P_v[x] = g_v  [P_{v1}[x], P_{v0}[x]]
    &= 	\begin{cases}
      g_v  & \text{if } v0 \text{ or }  v1 \text{ evaluates to } 0, \\
      g_v [ g_{v1},\,  g_{v0} ] & \text{if } v0 \text{ and }  v1 \text{ evaluate to } 1, \\
    \end{cases}\\
    &= 	\begin{cases}
      g_v  & \text{if } v \text{ evaluates to } 1, \\
      1 & \text{if } v \text{ evaluates to } 0.		
    \end{cases}
  \end{align*}
  Note that $[ g_{v1},\,  g_{v0}] = [ g_{v0},\,  g_{v1}]^{-1} = g_v^{-1}$ for the last equality.
  Thus, $P_v$ satisfies Equation (2). For $P_v^{-1}$ the analogous
  statement can be shown with the same calculation. 
  For a leaf $v \in \{0,1\}^d$, we have $\abs{g_v} \in 2^{\mathcal{O}(d)} = n^{\mathcal{O}(1)}$
  by Condition (a) from Definition 1 (recall that $d \in \mathcal{O}(\log n)$).
  Hence, $P_v^{-1}$ and $P_v$ have polynomial
  length in $n$. Finally, also $P_\eps$ has polynomial length in $n$ (with the same
  argument as for $g_\eps$; see the remark after Definition 1).

This gives us a non-uniform $\mathrm{AC}^0$-reduction (more precisely, a projection reduction) of $L(M)$ to $\mathrm{WP}(G)$. In order to obtain a DLOGTIME-reduction, we apply essentially the same construction. However, we need to introduce some padding so that the function mapping an index $i$ encoded in binary to the $i$-th instruction of $P_\varepsilon$ can be computed in DLINTIME. In order to show the last point, we need condition (d) of the uniform SENS definition (Definition 1).
\end{proof}

\section*{Direct proof for Thompson's group (Proposition 21)}
One can also prove the uniform SENS property for Thompson's group $F$ directly. Recall the infinite presentation $F=\langle x_0,x_1,x_2,\ldots\mid x_k^{x_i}=x_{k+1}\ (i<k)\rangle$.
\begin{authorclaim}{Proposition 21}
Let $g=x_3x_2^{-1}\in F$ and define $c_v$ for $v\in\{0,1\}^*$ inductively via
\[
c_\varepsilon=\varepsilon,\qquad c_{v0}=x_1c_v,\qquad c_{v1}=x_0^{-1}x_1c_v.
\]
Finally, for $d\in\mathbb N$ and $v\in\{0,1\}^{\le d}$ let $g_{d,v}=g^{c_v}$. Then $g_{d,v}=[g_{d,v0},g_{d,v1}]$ for all $d\in\mathbb N$ and $v\in\{0,1\}^{\le d}$ and $g_\varepsilon=g\ne1$ in $F$. In particular, $G$ is uniformly SENS.
\end{authorclaim}
\begin{proof}
Obviously, we have $g_\varepsilon=g$ in $F$. Moreover, since $x_2\ne x_3$ (which follows directly from the normal form theorem for $F$; see e.g.\ [12, Corollary 2.7]), we have $g\ne1$ in $F$ (Condition (c) of Definition 1). The identity $g=[g,g^{x_0^{-1}}]^{x_1}=[g^{x_1},g^{x_0^{-1}x_1}]$ is straightforward to check. Thus, we obtain Condition (b):
\[
g_{d,v}=g^{c_v}=[g^{x_1c_v},g^{x_0^{-1}x_1c_v}]
=[g^{c_{v0}},g^{c_{v1}}]=[g_{d,v0},g_{d,v1}].
\]
Moreover, since $|g_{d,v}|=|g|+2|c_v|\le|g|+4d$, we can pad with the identity symbol writing $g_{v,d}$ as a word of length $2^{\mu d}$ for some proper constant $\mu\in\mathbb N$ in order to meet Condition (a) (be aware that the lengths need to be computed over a finite standard generating set $\Sigma$, e.g.\ $\Sigma=\{1,x_0,x_0^{-1},x_1,x_1^{-1}\}$). This shows that $F$ is SENS.
Condition (d) of Definition 1 is also straightforward to check by introducing some more padding: we slightly change the definition of the $c_v$ by setting $c_\varepsilon=\varepsilon$, $c_{v0}=1x_1c_v$, and $c_{v1}=x_0^{-1}x_1c_v$ where $1\in\Sigma$. This new $c_v$ represents the same group element as the old $c_v$, but now we have $|c_v|=2|v|$ for all $v\in\{0,1\}^*$ and all letters at even positions are $x_1$, while letters at odd positions are either $1$ or $x_0^{-1}$ (the $(2j+1)$-st letter of $c_v$ depends on the $(|v|-j)$-th bit of $v$). Notice that we have $|g_{d,v}|=|g|+2|c_v|=|g|+4d$.
On input of $v\in\{0,1\}^d$, $i\in\mathbb N$ (encoded with $\mu d$ bits), and $a\in\Sigma$, first decide whether $1\le i\le2d$, or $2d<i\le2d+|g|$, or $2d+|g|<i\le4d+|g|$, or $i>4d+|g|$. This clearly can be done in DLINTIME (recall that $g$ is a constant). Now assume that $1\le i\le2d$ (i.e., $i$ points into the leading $c_v^{-1}$ of $g_{d,v}=c_v^{-1}gc_v$). If $i$ is odd, one accepts if and only if $a=x_1^{-1}$; if $i$ is even, one accepts if and only if the $i/2$-th bit of $v$ is zero and $a=1$ or if the $i/2$-th bit of $v$ is one and $a=x_0$. The other cases are similar.
\end{proof}

\section*{Direct proof for Grigorchuk's group}
\begin{authorclaim}{Proposition 22}
  Consider in the Grigorchuk group $G=\langle a,b,c,d\rangle$ the elements
  \[x = (abad)^2\quad \text{and} \quad y = x^b = babadabac.\]
  Define recursively elements $z_v\in\oneset{x,y,x^{-1},y^{-1}}$ for all
  $v\in\oneset{0,1}^*$ as follows:
  \begin{itemize}
  \item $z_\eps = x$;
  \item if $z_v$ is defined, then we define $z_{v0}$ and $z_{v1}$ according to the following table:
   \medskip
  \begin{center}
    \begin{tabular}{|c|c|c|}
      $z_v$ & $z_{v0}$ & $z_{v1}$ \\ \hline 
            $x$ & $x^{-1}$   & $y^{-1}$  \\ \hline
   $x^{-1}$ & $y^{-1}$ & $x^{-1}$  \\ \hline
          $y$ & $y$ & $x$ \\ \hline
  $y^{-1}$ & $x$ & $y$  
   \end{tabular}
 \end{center}
 \medskip
  \end{itemize}
  For every $d\in\mathbb N$  and $v \in \{0,1\}^{\le d}$ let 
  $g_{d,v} = z_v$ for $\abs v = d$ and $g_{d,v} = [g_{d,v0},g_{d,v1}]$ if $\abs v<d$.
  We then have $g_{d,\eps} \neq 1$ in $G$.
  In particular, $G$ is uniformly SENS.
\end{authorclaim}
 
\begin{proof}
  That $x\neq1 \neq y$ is easy to check by computing the action of $x$ and $y$ on the
  third level of the tree.
  Now the following equations are easy to check in $G$:
  \begin{align*}
    [x,y] &= \pair<1,{\pair<1,y^{-1}>}>,  \\ 
    [x^{-1},y^{-1}] &= \pair<1,{\pair<1,x>}>,  \\ 
    [y,x] &= \pair<1,{\pair<1,y>}>,  \\ 
    [y^{-1},x^{-1}] &= \pair<1,{\pair<1,x^{-1}>}>.  
  \end{align*}
  Hence, $[z_{v0},z_{v1}] = \pair<1,{\pair<1,z_v>}>$.
  The checks are tedious to compute by hand, but easy in the GAP
  package FR (note that vertices are numbered from $1$ in GAP and from
  $0$ here):
\begin{verbatim}
gap> LoadPackage("fr");
gap> AssignGeneratorVariables(GrigorchukGroup);
gap> x := (a*b*a*d)^2; y := x^b;
gap> Assert(0,Comm(x,y) = VertexElement([2,2],y^-1));
gap> Assert(0,Comm(x^-1,y^-1) = VertexElement([2,2],x));
gap> Assert(0,Comm(y,x) = VertexElement([2,2],y));
gap> Assert(0,Comm(y^-1,x^-1) = VertexElement([2,2],x^-1));
\end{verbatim}
  We claim that $g_{d,\eps} \neq 1$ in $G$. The equation  $[z_{v0},z_{v1}] = \pair<1,{\pair<1,z_v>}>$
  immediately implies that $g_{d,v}$ acts as $z_v$ on the subtree
  below vertex $1^{2(d-\abs v)}$
  and trivially elsewhere. In particular,
  $g_{d,\eps}$ acts as $z_\eps=x \neq 1$ on the subtree below vertex
  $1^{2d}$ and is non-trivial.
 	
  With this definition, the $g_{d,v}$ satisfy the definition of a SENS group. 
  Moreover, given some $v\in \oneset{0,1}^d$, $g_{d,v}$ can be
  computed in time $\Oh(d)$ by a deterministic finite state automaton with state set 
  $\oneset{x^{\pm1},y^{\pm1}}$.
\end{proof}


\begin{thebibliography}{51}
\bibitem{ref1} Ian Agol. The virtual Haken conjecture. Documenta Mathematica, 18:1045–1087, 2013. With an appendix by Ian Agol, Daniel Groves, and Jason Manning.
\bibitem{ref2} Sanjeev Arora and Boaz Barak. Computational Complexity - A Modern Approach. Cambridge University Press, 2009. doi:10.1017/CBO9780511804090.
\bibitem{ref3} David A. Mix Barrington. Bounded-width polynomial-size branching programs recognize exactly those languages in $\mathrm{NC}^1$ . Journal of Computer and System Sciences, 38(1):150–164, 1989. doi:10.1016/0022-0000(89)90037-8.
\bibitem{ref4} David A. Mix Barrington and Denis Thérien. Finite monoids and the fine structure of $\mathrm{NC}^1$ . Journal of the ACM, 35:941–952, 1988. doi:10.1145/48014.63138.
\bibitem{ref5} Laurent Bartholdi, Michael Figelius, Markus Lohrey, and Armin Weiß. Groups with ALOGTIME-hard word problems and PSPACE-complete compressed word problems. Technical report, arXiv.org, 2020. arXiv:1909.13781.
\bibitem{ref6} Laurent Bartholdi, Rostislav I. Grigorchuk, and Zoran Šuniḱ. Branch groups. In Handbook of algebra, Volume 3, pages 989–1112. Elsevier/North-Holland, Amsterdam, 2003. doi: 10.1016/S1570-7954(03)80078-5.
\bibitem{ref7} Laurent Bartholdi and Volodymyr V. Nekrashevych. Iterated monodromy groups of quadratic polynomials. I. Groups, Geometry, and Dynamics, 2(3):309–336, 2008. doi:10.4171/GGD/42.
\bibitem{ref8} Martin Beaudry, Pierre McKenzie, Pierre Péladeau, and Denis Thérien. Finite monoids: From word to circuit evaluation. SIAM Journal on Computing, 26(1):138–152, 1997. doi: 10.1137/S0097539793249530.
\bibitem{ref9} Richard Beigel and John Gill. Counting classes: Thresholds, parity, mods, and fewness. Theoretical Computer Science, 103(1):3–23, 1992. doi:10.1016/0304-3975(92)90084-S.
\bibitem{ref10} William W. Boone. The Word Problem. Annals of Mathematics, 70(2):207–265, 1959. URL: www.jstor.org/stable/1970103.
\bibitem{ref11} Daniel P. Bovet, Pierluigi Crescenzi, and Riccardo Silvestri. A uniform approach to define complexity classes. Theoretical Computer Science, 104(2):263–283, 1992. doi:10.1016/ 0304-3975(92)90125-Y.
\bibitem{ref12} John W. Cannon, William J. Floyd, and Walter R. Parry. Introductory notes on Richard Thompson’s groups. L’Enseignement Mathématique, 42(3):215–256, 1996.
\bibitem{ref13} Moses Charikar, Eric Lehman, Ding Liu, Rina Panigrahy, Manoj Prabhakaran, Amit Sahai, and Abhi Shelat. The smallest grammar problem. IEEE Transactions on Information Theory, 51(7):2554–2576, 2005. doi:10.1109/TIT.2005.850116.
\bibitem{ref14} Max Dehn. Über unendliche diskontinuierliche Gruppen. Mathematische Annalen, 71(1):116– 144, 1911. doi:10.1007/BF01456932.
\bibitem{ref15} Michael R. Garey and David S. Johnson. Computers and Intractability: A Guide to the Theory of NP–completeness. Freeman, 1979.
\bibitem{ref16} Max Garzon and Yechezkel Zalcstein. The complexity of Grigorchuk groups with application to cryptography. Theoretical Computer Science, 88(1):83–98, 1991. doi:10.1016/0304-3975(91) 90074-C.
\bibitem{ref17} Rostislav I. Grigorchuk. Burnside’s problem on periodic groups. Functional Analysis and Its Applications, 14:41–43, 1980. doi:10.1007/BF01078416.
\bibitem{ref18} Rostislav I. Grigorchuk and Zoran Šuniḱ. Asymptotic aspects of Schreier graphs and Hanoi Towers groups. C. R. Math. Acad. Sci. Paris, 342(8):545–550, 2006. doi:10.1016/j.crma. 2006.02.001.
\bibitem{ref19} Victor S. Guba and Mark V. Sapir. On subgroups of the R. Thompson group F and other diagram groups. Matematicheskii Sbornik, 190(8):3–60, 1999. doi:10.1070/ SM1999v190n08ABEH000419.
\bibitem{ref20} Narain Gupta and Saïd Sidki. On the Burnside problem for periodic groups. Mathematische Zeitschrift, 182(3):385–388, 1983. doi:10.1007/BF01179757.
\bibitem{ref21} Frédéric Haglund and Daniel T. Wise. Coxeter groups are virtually special. Advances in Mathematics, 224(5):1890–1903, 2010. doi:10.1016/j.aim.2010.01.011.
\bibitem{ref22} Ulrich Hertrampf. Über Komplexitätsklassen, die mit Hilfe von k-wertigen Funktionen definiert werden. Habilitationsschrift, Universität Würzburg, 1994.
\bibitem{ref23} Ulrich Hertrampf. The shapes of trees. In Proceedings of COCOON 1997, volume 1276 of Lecture Notes in Computer Science, pages 412–421. Springer, 1997. doi:10.1007/BFb0045108.
\bibitem{ref24} Ulrich Hertrampf. Algebraic acceptance mechanisms for polynomial time machines. SIGACT News, 31(2):22–33, 2000. doi:10.1145/348210.348215.
\bibitem{ref25} Ulrich Hertrampf, Clemens Lautemann, Thomas Schwentick, Heribert Vollmer, and Klaus W. Wagner. On the power of polynomial time bit-reductions. In Proceedings of the Eighth Annual Structure in Complexity Theory Conference, pages 200–207. IEEE Computer Society Press, 1993. doi:10.1109/SCT.1993.336526.
\bibitem{ref26} Ulrich Hertrampf, Heribert Vollmer, and Klaus Wagner. On balanced versus unbalanced computation trees. Mathematical Systems Theory, 29(4):411–421, 1996. doi:10.1007/BF01192696.
\bibitem{ref27} Derek F. Holt, Markus Lohrey, and Saul Schleimer. Compressed decision problems in hyperbolic groups. In Proceedings of STACS 2019, volume 126 of LIPIcs, pages 37:1–37:16. Schloss Dagstuhl - Leibniz-Zentrum für Informatik, 2019. URL: http://www.dagstuhl.de/dagpub/ 978-3-95977-100-9.
\bibitem{ref28} Derek F. Holt, Sarah Rees, and Claas E. Röver. Groups, Languages and Automata, volume 88 of London Mathematical Society Student Texts. Cambridge University Press, 2017. doi: 10.1017/9781316588246.
\bibitem{ref29} Birgit Jenner, Pierre McKenzie, and Denis Thérien. Logspace and logtime leaf languages. Information and Computation, 129(1):21–33, 1996. doi:10.1006/inco.1996.0071.
\bibitem{ref30} Howard J. Karloff and Walter L. Ruzzo. The iterated mod problem. Information and Computation, 80(3):193–204, 1989. doi:10.1016/0890-5401(89)90008-4.
\bibitem{ref31} Daniel König and Markus Lohrey. Evaluation of circuits over nilpotent and polycyclic groups. Algorithmica, 80(5):1459–1492, 2018. doi:10.1007/s00453-017-0343-z.
\bibitem{ref32} Daniel König and Markus Lohrey. Parallel identity testing for skew circuits with big powers and applications. International Journal of Algebra and Computation, 28(6):979–1004, 2018. doi:10.1142/S0218196718500431.
\bibitem{ref33} Jörg Lehnert and Pascal Schweitzer. The co-word problem for the Higman-Thompson group is context-free. Bulletin of the London Mathematical Society, 39(2):235–241, February 2007. doi:10.1112/blms/bdl043.
\bibitem{ref34} Martin W. Liebeck, Eamonn A. O’Brien, Aner Shalev, and Pham Huu Tiep. The Ore conjecture. Journal of the European Mathematical Society, 12(4):939–1008, 2010. doi:10.4171/JEMS/220.
\bibitem{ref35} Yury Lifshits and Markus Lohrey. Querying and embedding compressed texts. In Proceedings of MFCS 2006, volume 4162 of Lecture Notes in Computer Science, pages 681–692. Springer, 2006. doi:10.1007/11821069\_59.
\bibitem{ref36} Richard J. Lipton and Yechezkel Zalcstein. Word problems solvable in logspace. Journal of the Association for Computing Machinery, 24(3):522–526, 1977. doi:10.1145/322017.322031.
\bibitem{ref37} Markus Lohrey. Leaf languages and string compression. Information and Computation, 209(6):951–965, 2011. doi:10.1016/j.ic.2011.01.009.
\bibitem{ref38} Markus Lohrey. The Compressed Word Problem for Groups. Springer Briefs in Mathematics. Springer, 2014. doi:10.1007/978-1-4939-0748-9.
\bibitem{ref39} Volodymyr Nekrashevych. Self-similar groups, volume 117 of Mathematical Surveys and Monographs. American Mathematical Society, Providence, RI, 2005. doi:10.1090/surv/117.
\bibitem{ref40} Piotr S. Novikov. On the algorithmic unsolvability of the word problem in group theory. Trudy Mat. Inst. Steklov, pages 1–143, 1955. In Russian. URL: http://mi.mathnet.ru/eng/tm1180.
\bibitem{ref41} David Robinson. Parallel Algorithms for Group Word Problems. PhD thesis, University of California, San Diego, 1993.
\bibitem{ref42} Joseph J. Rotman. An Introduction to the Theory of Groups (fourth edition). Springer, 1995.
\bibitem{ref43} Amir Shpilka and Amir Yehudayoff. Arithmetic circuits: A survey of recent results and open questions. Foundations and Trends in Theoretical Computer Science, 5(3-4):207–388, 2010. doi:10.1561/0400000039.
\bibitem{ref44} Hans-Ulrich Simon. Word problems for groups and contextfree recognition. In Proceedings of FCT 1979, pages 417–422. Akademie-Verlag, 1979.
\bibitem{ref45} Roman Smolensky. Algebraic methods in the theory of lower bounds for boolean circuit complexity. In Proceedings of STOC 1987, pages 77–82. ACM, 1987. doi:10.1145/28395. 28404.
\bibitem{ref46} Jacques Tits. Free subgroups in linear groups. Journal of Algebra, 20(2):250–270, 1972. doi:10.1016/0021-8693(72)90058-0.
\bibitem{ref47} Heribert Vollmer. Introduction to Circuit Complexity. Springer, Berlin, 1999. doi:10.1007/ 978-3-662-03927-4.
\bibitem{ref48} Stephan Waack. The parallel complexity of some constructions in combinatorial group theory. Journal of Information Processing and Cybernetics EIK, 26:265–281, 1990. doi: 10.1007/BFb0029647.
\bibitem{ref49} Jan Philipp Wächter and Armin Weiß. An automaton group with pspace-complete word problem. In Proceedings of STACS 2020, volume 154 of LIPIcs, pages 6:1–6:17. Schloss Dagstuhl - Leibniz-Zentrum für Informatik, 2020. URL: https://www.dagstuhl.de/dagpub/ 978-3-95977-140-5.
\bibitem{ref50} John S. Wilson. Embedding theorems for residually finite groups. Mathematische Zeitschrift, 174(2):149–157, 1980. doi:10.1007/BF01293535.
\bibitem{ref51} Daniel T. Wise. Research announcement: the structure of groups with a quasiconvex hierarchy. Electronic Research Announcements in Mathematical Sciences, 16:44–55, 2009. URL: http: //aimsciences.org//article/id/8d3fc128-8d02-4fee-af67-38001ca1d0ac.
\end{thebibliography}
\end{document}
